\section{Conclusion}
\label{s-conclusion}

The fixed-\(\omega\) problem has complementary one-dimensional
descriptions. A scalar equation gives the relaxed optimizer, while
ordering increments by Nicolas's score gives every integer optimizer as a
score-ordered prefix and supplies a finite computation. The integer exponents
are within one step of the relaxed exponents. Transferring their rounding
loss to a prime integral gives
\(2\sqrt2/(\sqrt{p_m}\log p_m)\), with equal contributions from the two
sides of the second-exponent switch.

Classical CA shells locate the same switch and explain the common
coefficient in the Axler--Nicolas size-constrained comparison, although
the feasible sets and the second terms differ. Combining the unconditional
gap with Nicolas's Mertens formula gives the eventual fixed-\(\omega\)
criterion. The same score ordering also characterizes the bounded-support
problem \(\omega(n)\leq m\).
